\section{Results}

\subsection{A literature-linked human TE graph with explicit provenance}

The current Neo4j snapshot contains 225 TE nodes and 2,308 Paper nodes. It also contains 12,444 directed \texttt{BIO\_RELATION} relationships connecting TEs to biomedical entities. Every current biological relation has a normalized predicate, at least one PMID field and a supporting-PMID count. These properties allow a displayed edge to be traced to the literature record used to support it, although provenance does not convert an extracted association into causal evidence.

The biomedical schema contains 11 entity categories. The most numerous current labels are Function (3,683 nodes), Gene (1,280), Protein (1,089), Disease (676), RNA (588), Mutation (377) and Pharmaceutical (293). Smaller categories include Toxin (67), Lipid (26), Peptide (23) and Carbohydrate (12). DiseaseCategory (767), Paper (2,308) and taxonomy-related nodes are represented separately because they serve organizational or provenance roles rather than the 11 biomedical entity categories.

\begin{table*}[!t]
\caption{Current TE-KG content snapshot. Counts describe different runtime units and must not be summed. The snapshot was verified on 31 July 2026.\label{tab:snapshot}}
\begin{tabularx}{\textwidth}{@{}l r X X@{}}
\toprule
Component and unit & Count & Definition & Source and boundary \\
\midrule
Neo4j TE nodes & 225 & TE entities in the \texttt{tekg3} graph & Live label-count query; not identical to Browse entries \\
Neo4j Paper nodes & 2,308 & Retained publications represented in the graph & Live label-count query; agrees with the final literature corpus \\
Directed biological relations & 12,444 & Stored \texttt{BIO\_RELATION} relationships & Counted once in the stored direction; all current records contain predicate and PMID provenance fields \\
Biomedical entity categories & 11 & Carbohydrate, Disease, Function, Gene, Lipid, Mutation, Peptide, Pharmaceutical, Protein, RNA and Toxin & Excludes Paper, DiseaseCategory and taxonomy labels \\
Browse catalogue entries & 276 & Versioned TE catalogue records & MySQL-backed Browse API; a different record unit from graph TE nodes \\
Taxonomy-classified TEs & 192 of 225 & TE nodes with a taxonomy class in the live summary & Neo4j-backed taxonomy API \\
Expression samples & 1,158 & 205 normal tissue, 307 normal primary cell and 646 cancer cell line & Three source groups; not a matched cross-context cohort \\
Searchable co-expression entries & 784 & 285 TE and 499 Gene entries & Approved display catalogue across three contexts; not the number of full offline-network features \\
Download files & 10 & Six expression, two graph and two taxonomy files & Current web route; public archival release still required \\
\botrule
\end{tabularx}
\end{table*}



\subsection{Classification and catalogue views retain distinct record units}

The live taxonomy summary contains 225 TEs, of which 192 have an assigned taxonomy class. The summary also reports 188 standard leaves and separates a 36-record all-source tree from a 189-record RMSK plus RepBase tree. These values describe API-defined taxonomy views rather than independent biological totals. The Browse catalogue contains 276 versioned entries. This larger Browse count is reported separately from the Neo4j TE-node count because the two components expose different record units and should not be combined into a single TE total.

The classification interface offers both a hierarchical tree and a force-directed graph while using the same API-backed taxonomy data. Its legend can temporarily hide or restore TE classes without changing the underlying result set. This interaction supports inspection of dense classifications, but it does not introduce a second taxonomy source or modify classification membership.

\subsection{Expression and co-expression add three context classes}

The active expression matrices contain 205 normal-tissue, 307 normal-primary-cell and 646 cancer-cell-line samples. These 1,158 samples provide three context classes for TE and gene expression inspection. The resource keeps the contexts separate because they originate from different studies and sampling designs. The current data do not support treating them as a matched experimental comparison.

The production co-expression catalogue spans the same three context classes and includes approved searchable entries for 285 TEs and 499 genes. Across contexts, 849 display recommendations are available: 17 are marked as core cases, 287 as high-confidence entries, 542 as searchable entries and three as not recommended for default display. These tiers govern presentation and searchability; they are not statistical significance categories. All displayed edges already satisfy the fixed correlation and FDR thresholds, and a returned display network is a bounded view of a larger offline result.

\subsection{Connected interfaces expose complementary evidence routes}

TE-KG provides several entry routes suited to different questions. Browse and Search resolve TE names and lead to structured identity records. Knowledge Graph displays literature-linked biomedical relations and exposes PubMed evidence. Path searches for a connection between selected entities. Expression shows context-specific abundance profiles, and Co-expression displays approved TE-gene correlation neighbourhoods. Download currently exposes ten files comprising six expression artifacts, two graph artifacts and two taxonomy artifacts.

These routes are connected at the level of TE identity but retain separate evidence semantics. For example, opening a disease-linked graph edge gives access to the supporting paper rather than a claim that the TE causes the disease. Likewise, opening a co-expression neighbourhood reports the context and correlation threshold rather than a regulatory mechanism. This separation is part of the user-visible design, not only an implementation detail.

\subsection{Natural-language follow-up remains bounded to the current session}

Agent and DeepThink can retrieve several local evidence layers in response to an English question and can interpret a follow-up turn within the current conversation. PubMed identifiers in generated answers are rendered as links to PubMed. The writing layer translates evidence limitations into ordinary scientific language instead of displaying internal routing, confidence labels or other developer-facing diagnostics. For example, an answer can describe a relation as an association that requires experimental validation without asking the reader to interpret an internal status code.

The current evaluation supports functional coverage but not a global accuracy claim. A fixed 13-question baseline exercised all 12 retrieval roles, and subsequent regression checks protected graph-routing, literature-link and multi-plugin behavior after targeted fixes. A larger historical 36-question run included failures that were later corrected, but the entire suite has not been rerun as one version-locked benchmark. We therefore report this interface as a bounded retrieval route and retain individual known limitations rather than publishing the historical pass fraction as current performance.

\subsection{Cross-layer L1HS use case}

L1HS was selected as the principal worked example because it can be resolved across the current classification, graph, sequence, genome, expression and co-expression interfaces. The final figure will preserve the distinct result type in each panel and link all literature statements to verified PMIDs. \draftnote{Freeze the exact L1HS records, expression source table, co-expression neighbourhood and one or two PMID-supported graph relations before this subsection is considered manuscript-ready. Do not fill missing layers from general model knowledge.}
